% ============================================================
%  CUADERNILLO DE ESTUDIO — CLASE 15
%  Seguridad de la Información y de Sistemas
%  Lic. en Gestión de Negocios Digitales — UCA
%  Compilar con: pdflatex (dos pasadas para el índice)
% ============================================================

\documentclass[12pt,a4paper]{article}

% ---------- Paquetes ----------
\usepackage[utf8]{inputenc}
\usepackage[spanish]{babel}
\usepackage[T1]{fontenc}
\usepackage{lmodern}
\usepackage[margin=2.2cm,headheight=14pt]{geometry}
\usepackage{amsmath}
\usepackage{amssymb}
\usepackage{enumitem}
\usepackage{booktabs}
\usepackage{array}
\usepackage{tabularx}
\usepackage[table]{xcolor}
\usepackage{tcolorbox}
\usepackage{graphicx}
\usepackage{fancyhdr}
\usepackage{titlesec}
\usepackage{hyperref}
\usepackage{multicol}
\usepackage{pifont}
\usepackage{wasysym}
\usepackage{tikz}
\usetikzlibrary{positioning,arrows.meta,shapes.geometric}

% ---------- Colores ----------
\definecolor{azulUCA}{HTML}{003366}
\definecolor{doradoUCA}{HTML}{B8860B}
\definecolor{grisClaro}{HTML}{F2F2F2}
\definecolor{rojoAlerta}{HTML}{C0392B}
\definecolor{verdeOK}{HTML}{27AE60}
\definecolor{azulCielo}{HTML}{2980B9}
\definecolor{violetaIA}{HTML}{8E44AD}
\definecolor{doradoFinal}{HTML}{7D6608}
\definecolor{azulExamen}{HTML}{154360}

% ---------- tcolorbox estilos ----------
\tcbuselibrary{skins,breakable}

\newtcolorbox{cajaTitulo}[1][]{
  colback=azulUCA, colframe=azulUCA, coltext=white,
  fonttitle=\Large\bfseries, arc=4mm, #1
}

\newtcolorbox{cajaDefensa}[1][]{
  colback=doradoFinal!10, colframe=doradoFinal,
  fonttitle=\bfseries\color{doradoFinal},
  title={\ding{73}\; Defensa Oral}, arc=3mm, breakable, #1
}

\newtcolorbox{cajaExamen}[1][]{
  colback=azulExamen!8, colframe=azulExamen,
  fonttitle=\bfseries\color{azulExamen},
  title={\ding{115}\; Evaluaci\'on Final}, arc=3mm, breakable, #1
}

\newtcolorbox{cajaActividad}[1][]{
  colback=grisClaro, colframe=azulUCA,
  fonttitle=\bfseries\color{azulUCA},
  title={\ding{45}\; Actividad}, arc=3mm, breakable, #1
}

\newtcolorbox{cajaNota}[1][]{
  colback=doradoUCA!10, colframe=doradoUCA,
  fonttitle=\bfseries, title=Concepto clave, arc=3mm, breakable, #1
}

\newtcolorbox{cajaCierre}[1][]{
  colback=verdeOK!10, colframe=verdeOK,
  fonttitle=\bfseries\color{verdeOK}, title=Cierre del Curso, arc=3mm, breakable, #1
}

% ---------- Header / Footer ----------
\pagestyle{fancy}
\fancyhf{}
\fancyhead[L]{\small\textcolor{azulUCA}{Seguridad de la Informaci\'on y de Sistemas}}
\fancyhead[R]{\small\textcolor{azulUCA}{Cuadernillo --- Clase 15 (Defensa Final)}}
\fancyfoot[C]{\thepage}
\fancyfoot[L]{\small\textcolor{gray}{UCA}}
\renewcommand{\headrulewidth}{0.4pt}
\renewcommand{\footrulewidth}{0.4pt}

% ---------- Títulos de sección ----------
\titleformat{\section}
  {\color{azulUCA}\Large\bfseries}
  {\thesection.}{0.6em}{}
\titleformat{\subsection}
  {\color{azulCielo}\large\bfseries}
  {\thesubsection}{0.6em}{}

% ---------- Enlaces ----------
\hypersetup{colorlinks=true, linkcolor=azulUCA, urlcolor=azulCielo}

% ---------- Checkbox ----------
\newcommand{\checkboxVacio}{$\square$\;}
\newcommand{\checkMark}{\ding{51}}

% ---------- Reducir padding de tablas ----------
\setlength{\tabcolsep}{5pt}

% ============================================================
\begin{document}

% ============================================================
%  PORTADA
% ============================================================
\begin{titlepage}
\centering
\vspace*{1cm}

\begin{cajaTitulo}[width=\textwidth, halign=center]
  SEGURIDAD DE LA INFORMACI\'ON Y DE SISTEMAS\\[4pt]
  {\large Licenciatura en Gesti\'on de Negocios Digitales --- UCA}
\end{cajaTitulo}

\vspace{1.5cm}

{\Huge\bfseries\color{azulUCA} CUADERNILLO DE ESTUDIO}\\[8pt]
{\LARGE\color{doradoUCA} CLASE 15}\\[12pt]
{\Large\itshape ``Presentaci\'on y defensa del\\[2pt] Trabajo Pr\'actico Integrador''}\\[6pt]
{\large\color{doradoFinal}\ding{73}\; El d\'ia de la verdad}\\[1.5cm]

\begin{tabular}{rl}
\textbf{Profesor:}   & Mg.\ Ing.\ Rodrigo Atilio Elgueta \\[4pt]
\textbf{Eje:}        & Cierre de curso y Evaluaci\'on Final \\[4pt]
\end{tabular}

\vfill

\begin{tcolorbox}[colback=grisClaro, colframe=azulUCA, arc=3mm, width=0.9\textwidth]
  \centering
  {\bfseries Datos del/la estudiante}\\[10pt]
  Nombre y apellido: \underline{\hspace{5cm}}\\[10pt]
  Grupo N.\textdegree: \underline{\hspace{5cm}}\\[10pt]
  Negocio Digital (TP): \underline{\hspace{5cm}}\\[10pt]
  Fecha: \underline{\hspace{5cm}}
\end{tcolorbox}

\vspace{1cm}
{\small\textcolor{gray}{Material de uso exclusivo para la cursada.}}
\end{titlepage}

% ============================================================
%  ÍNDICE
% ============================================================
\tableofcontents
\newpage

% ============================================================
%  SECCIÓN 1 — INTRODUCCIÓN
% ============================================================
\section{El gran d\'ia: Defensa del TP Integrador}

Bienvenidos a la \'ultima clase del cuatrimestre. Hoy no hay exposici\'on te\'orica del docente, ni talleres guiados. Hoy los protagonistas son ustedes. 

\begin{cajaExamen}[title={Esta clase reemplaza el examen final}]
Seg\'un el programa de la materia, la defensa oral de hoy \textbf{reemplaza la instancia tradicional de examen final}.\\[4pt]
\textbf{Formato:}
\begin{itemize}[leftmargin=2em]
    \item \textbf{10 minutos} de exposici\'on grupal ininterrumpida.
    \item \textbf{5 minutos} de preguntas individuales cruzadas por parte del docente (y el evaluador invitado, si lo hay).
\end{itemize}
\textbf{Nota:} 1 a 10. Se aprueba con 4.
\end{cajaExamen}

\begin{cajaNota}[title={El cambio de rol}]
Durante 14 clases, el docente fue su gu\'ia. Hoy, el docente (y la audiencia) son \textbf{clientes, directivos o auditores}. Ustedes son los consultores senior de ciberseguridad que vienen a presentar el ``Programa de Seguridad'' para el negocio digital que eligieron. V\'endanlo, defi\'endanlo y demuestren que dominan el tablero completo.
\end{cajaNota}

\newpage
% ============================================================
%  SECCIÓN 2 — LA RÚBRICA OFICIAL
% ============================================================
\section{La R\'ubrica Oficial de Evaluaci\'on}

Esta es la r\'ubrica con la que ser\'an calificados hoy. T\'enganla a la vista mientras preparan su exposici\'on y mientras escuchan a sus compa\~neros.

\renewcommand{\arraystretch}{1.5}
\begin{tabularx}{\textwidth}{|X|>{\centering\arraybackslash}p{1.5cm}|}
\hline
\rowcolor{doradoFinal!20}
\textbf{Dimensi\'on evaluada} & \textbf{Peso} \\
\hline
\textbf{1. Comprensi\'on conceptual}\\
Aplican correctamente los pilares CID, la gesti\'on de riesgos, el marco legal y el gobierno de la seguridad al caso elegido. Demuestran que entienden el ``por qu\'e'' de cada decisi\'on. & \textbf{25\%} \\
\hline
\textbf{2. Viabilidad y aplicaci\'on pr\'actica}\\
Las propuestas (controles, PSA, plan de incidentes) son concretas, realistas y adecuadas al tama\~no y contexto del negocio elegido (no proponen controles de multinacional para una pyme). & \textbf{25\%} \\
\hline
\textbf{3. Pensamiento cr\'itico sobre el uso de IA}\\
El Anexo de IA documenta con honestidad qu\'e se us\'o, y muestra evidencia real de verificaci\'on y correcci\'on humana (no es una copia acr\'itica de la salida de la m\'aquina). & \textbf{20\%} \\
\hline
\textbf{4. Trabajo en equipo y proceso}\\
Evidencia de avances consistentes a lo largo del cuatrimestre y participaci\'on equilibrada de todos los integrantes del grupo. & \textbf{15\%} \\
\hline
\textbf{5. Comunicaci\'on y defensa oral}\\
Claridad expositiva, manejo del tiempo y calidad de las respuestas individuales durante la sesi\'on de preguntas. & \textbf{15\%} \\
\hline
\rowcolor{grisClaro}
\textbf{TOTAL} & \textbf{100\%} \\
\hline
\end{tabularx}

\newpage
% ============================================================
%  SECCIÓN 3 — GUÍA PARA LA EXPOSICIÓN (10 MIN)
% ============================================================
\section{Gu\'ia para los 10 minutos de exposici\'on}

Diez minutos pasan volando. No intenten leer el dossier de 15 p\'aginas; tienen que \textbf{contar la historia} de c\'omo protegieron su negocio.

\begin{cajaDefensa}[title={Estructura sugerida (Cronometrada)}]
\renewcommand{\arraystretch}{1.4}
\small
\begin{tabularx}{\textwidth}{|>{\bfseries\color{doradoFinal}}p{2.2cm}|X|>{\centering\arraybackslash}p{1.2cm}|}
\hline
\rowcolor{grisClaro}
\textbf{Bloque} & \textbf{Qu\'e mostrar / Qu\'e decir} & \textbf{Tiempo} \\
\hline
1. El Negocio & Presentar la empresa, su modelo de negocio y sus activos m\'as cr\'iticos. & 1.5 min \\
\hline
2. El Riesgo & Mostrar el ``Mapa de Calor'' (Matriz P$\times$I). ¿Cu\'ales son los 2 o 3 riesgos que podr\'ian destruir el negocio? & 2 min \\
\hline
3. La Soluci\'on & Los controles priorizados (t\'ecnicos, humanos, legales). Conectar el BIA con la Continuidad. & 2.5 min \\
\hline
4. Gobierno y Ley & C\'omo se ordenan internamente (PSA) y c\'omo respetan al usuario (Privacidad / ARCO). & 2 min \\
\hline
5. El factor IA & Qu\'e aprendieron al auditar la IA generativa en su Plan de Incidentes. & 1 min \\
\hline
6. Cierre & La conclusi\'on ejecutiva: ¿Por qu\'e este negocio ahora es confiable? & 1 min \\
\hline
\rowcolor{grisClaro}
\textbf{TOTAL} & & \textbf{10 min} \\
\hline
\end{tabularx}
\end{cajaDefensa}

\begin{cajaNota}[title={Consejos de oro para la presentaci\'on}]
\begin{itemize}[leftmargin=2em]
    \item \textbf{Hablen en idioma ``negocio'', no solo en idioma ``t\'ecnico''}: al directivo le importa el impacto econ\'omico y reputacional, no solo el algoritmo de cifrado.
    \item \textbf{No lean diapositivas}: usen esquemas visuales, matrices de calor, diagramas.
    \item \textbf{Repartan el tiempo equitativamente}: todos los integrantes deben hablar.
    \item \textbf{Si se pasan de tiempo, el docente los cortar\'a}: practiquen con cron\'ometro.
\end{itemize}
\end{cajaNota}

\newpage
% ============================================================
%  SECCIÓN 4 — GUÍA PARA LAS PREGUNTAS (5 MIN)
% ============================================================
\section{Gu\'ia para los 5 minutos de preguntas individuales}

Terminada la exposici\'on, el docente (y el jurado) har\'a preguntas \textbf{dirigidas a integrantes espec\'ificos}. El objetivo no es ``tomar examen de memoria'', sino verificar que \textbf{todos} participaron y comprenden la l\'ogica detr\'as del trabajo.

\subsection{¿Qu\'e se eval\'ua en esta instancia?}
\begin{itemize}[leftmargin=2em]
    \item \textbf{Agilidad mental:} ¿Pueden justificar una decisi\'on sobre la marcha?
    \item \textbf{Honestidad intelectual:} Si no saben algo, es mejor decir ``Ese punto lo investig\'o mi compa\~nero, pero mi lectura es...'' o ``No lo contemplamos, pero hoy lo resolver\'ia as\'i...''.
    \item \textbf{Criterio propio:} Especialmente en las preguntas sobre el uso de IA y el marco legal.
\end{itemize}

\subsection{Preguntas ``comod\'in'' que pueden aparecer}
\begin{cajaExamen}[title={Prep\'arense para estas preguntas}]
\begin{itemize}[leftmargin=2em]
    \item \emph{``Si ma\~nana tienen un presupuesto de solo \$1.000 para seguridad, ¿en qu\'e lo gastan y por qu\'e?''}
    \item \emph{``Un cliente ejerce su derecho de supresi\'on (ARCO), pero nos pide que borremos sus facturas de los \'ultimos 5 a\~nos. ¿Qu\'e le responden?''}
    \item \emph{``En su Anexo de IA, corrigieron una recomendaci\'on de la m\'aquina. ¿Por qu\'e la IA se equivoc\'o en ese contexto?''}
    \item \emph{``¿Qu\'e pasa si el proveedor de nube donde est\'an alojados sufre un ransomware? ¿Qu\'e dice su BIA al respecto?''}
    \item \emph{``¿C\'omo convencen al CEO de que firmar la Pol\'itica de Seguridad Aceptable (PSA) no es solo un tr\'amite burocr\'atico?''}
\end{itemize}
\end{cajaExamen}

\newpage
% ============================================================
%  SECCIÓN 5 — FICHA DE EVALUACIÓN (DOCENTE / JURADO)
% ============================================================
\section{Ficha de Evaluaci\'on (Uso del Docente / Jurado)}

\begin{cajaNota}[title={Para la c\'atedra}]
Esta ficha puede ser utilizada por el docente o por el evaluador invitado para calificar a cada grupo en tiempo real durante la defensa.
\end{cajaNota}

\vspace{6pt}
\noindent\textbf{Grupo / Empresa Evaluada:} \hrulefill\\[6pt]
\noindent\textbf{Integrantes presentes:} \hrulefill

\vspace{8pt}
\renewcommand{\arraystretch}{1.5}
\begin{tabularx}{\textwidth}{|X|>{\centering\arraybackslash}p{2.5cm}|}
\hline
\rowcolor{azulExamen!15}
\textbf{Dimensi\'on} & \textbf{Puntaje (0 a 10)} \\
\hline
\textbf{1. Comprensi\'on conceptual} (CID, riesgos, ley, gobierno) & \hrulefill \\
\hline
\textbf{2. Viabilidad pr\'actica} (realismo para el tama\~no del negocio) & \hrulefill \\
\hline
\textbf{3. Pensamiento cr\'itico IA} (Anexo de IA, auditor\'ia humana) & \hrulefill \\
\hline
\textbf{4. Trabajo en equipo} (participaci\'on equilibrada, proceso) & \hrulefill \\
\hline
\textbf{5. Comunicaci\'on y defensa} (exposici\'on + preguntas individuales) & \hrulefill \\
\hline
\rowcolor{grisClaro}
\textbf{PROMEDIO FINAL (Nota 1-10)} & \textbf{\hrulefill} \\
\hline
\end{tabularx}

\vspace{8pt}
\begin{tabularx}{\textwidth}{|X|}
\hline
\textbf{Comentarios / Devoluci\'on cualitativa principal:} \\[4cm]
\hline
\end{tabularx}

\newpage
% ============================================================
%  SECCIÓN 6 — ENCUESTA DE SATISFACCIÓN FINAL
% ============================================================
\section{Encuesta de satisfacci\'on final}

El programa de esta materia fue dise\~no en base a una encuesta de expectativas que ustedes respondieron en la Clase 1 (donde pidieron m\'as nube, hacking \'etico, IA y menos examen tradicional). Hoy, al cerrar el c\'irculo, necesitamos su feedback honesto para mejorar la materia para las pr\'oximas cohortes.

\begin{cajaActividad}[title={Reflexi\'on individual y an\'onima}]
Respondan con sinceridad. El docente leer\'a estas encuestas una vez cargadas las notas finales.
\end{cajaActividad}

\vspace{4pt}
\renewcommand{\arraystretch}{1.6}
\begin{tabularx}{\textwidth}{|X|}
\hline
\textbf{1. Sobre los contenidos: ¿Qu\'e tema de los que pedimos en la encuesta inicial (nube, IA, hacking, cookies) les pareci\'o m\'as \'util para su perfil de negocios digitales?} \\[2.5cm]
\hline
\textbf{2. Sobre la metodolog\'ia: Reemplazamos el examen final tradicional por este TP Integrador con defensa oral. ¿Sienten que evalu\'o mejor sus conocimientos? ¿Qu\'e le cambiar\'ian?} \\[3cm]
\hline
\textbf{3. Sobre el uso de IA: ¿C\'omo cambi\'o su forma de ver la Inteligencia Artificial al tener que ``auditarla'' en lugar de solo usarla como un buscador glorificado?} \\[3cm]
\hline
\textbf{4. El momento ``Aha!'': ¿En qu\'e clase o actividad espec\'ifica hicieron ``clic'' y entendieron por qu\'e la seguridad es un habilitador de negocios y no solo un problema de inform\'aticos?} \\[2.5cm]
\hline
\end{tabularx}

\newpage
% ============================================================
%  SECCIÓN 7 — CIERRE Y DESPEDIDA
% ============================================================
\section{Reflexi\'on Final: De la Clase 1 a hoy}

\begin{cajaCierre}[title={El c\'irculo se cierra}]
Hagamos un viaje en el tiempo de apenas 15 semanas.\\[6pt]
En la \textbf{Clase 1}, entraron a un link desde sus celulares. Una p\'agina aparentemente inocua accedi\'o a la c\'amara de sus dispositivos y los expuso en un panel de administraci\'on. En ese momento, la mayor\'ia sinti\'o incomodidad, sorpresa o miedo. Eran \textbf{usuarios vulnerables} en un mundo digital que no terminaban de entender.\\[6pt]
Hoy, en la \textbf{Clase 15}, ustedes son los que dise\~nan la arquitectura de ese negocio.\\[4pt]
\ding{51}\; Ustedes saben \textbf{por qu\'e} esa c\'amara se activ\'o (falta de consentimiento, violaci\'on del pilar de Confidencialidad).\\
\ding{51}\; Ustedes saben \textbf{qu\'e riesgo} corre la empresa que hizo eso (multas de la Ley 25.326, p\'erdida de reputaci\'on).\\
\ding{51}\; Ustedes saben \textbf{c\'omo prevenirlo} (Privacy by Design, cifrado, pol\'iticas de IAM, auditor\'ias).\\
\ding{51}\; Y ustedes saben \textbf{c\'omo responder} si, a\'un con todas las defensas, un incidente ocurre (Junta de Crisis, BIA, Plan de Continuidad).
\end{cajaCierre}

\vspace{6pt}

\begin{center}
\begin{tikzpicture}[
  box/.style={rectangle, draw=azulUCA, fill=azulUCA!10,
              text width=4.5cm, align=center, rounded corners=4pt,
              font=\small\bfseries, minimum height=2cm},
  arr/.style={-{Stealth[length=4mm]}, thick, doradoFinal, line width=1.5pt}
]
  \node[box, fill=rojoAlerta!15, draw=rojoAlerta, text width=5cm] (c1) 
    {CLASE 1\\``¿C\'omo me hackearon la c\'amara sin que me d\'e cuenta?''\\{\scriptsize (El usuario vulnerable)}};
  
  \node[box, fill=verdeOK!20, draw=verdeOK, text width=5cm, right=3cm of c1] (c15) 
    {CLASE 15\\``Este es nuestro Programa de Seguridad, nuestros Controles y nuestra Gobernanza.''\\{\scriptsize (El profesional de negocios digitales)}};

  \draw[arr] (c1) -- node[above, font=\small\itshape, text=doradoFinal] {15 semanas de cursada} (c15);
\end{tikzpicture}
\end{center}

\vspace{10pt}

\begin{cajaNota}[title={Palabras finales del docente}]
La seguridad de la informaci\'on ya no es una ``caja negra'' delegada al departamento de IT. Para ustedes, futuros l\'ideres de negocios digitales, la seguridad es \textbf{confianza}, es \textbf{supervivencia} y es \textbf{ventaja competitiva}.\\[6pt]
El tablero est\'a completo. Tienen los fundamentos, la matriz de riesgos, la ley, la nube, la criptograf\'ia y el gobierno. Ahora, salgan a construir negocios digitales que sean no solo rentables, sino \textbf{robustos, \'eticos y confiables}.\\[6pt]
¡Muchas gracias por el cuatrimestre y mucho \'exito en lo que viene!
\end{cajaNota}

\vfill
\begin{center}
\Large\bfseries\color{azulUCA}
¡FELICITACIONES!\\[6pt]
\normalsize\color{gray}
Han completado Seguridad de la Informaci\'on y de Sistemas.\\
Licenciatura en Gesti\'on de Negocios Digitales --- UCA.
\end{center}

\end{document}